La evidencia sugiere que los grafos Mapper seleccionados son utiles para priorizar regiones de inspeccion, siempre que se separen cuatro capas: senal fisica u orbital, sesgo observacional, dependencia de imputacion y sensibilidad al lens. El grafo orbital \texttt{orbital\_pca2\_cubes10\_overlap0p35} es un candidato importante por su complejidad y baja imputacion, pero la auditoria de sesgo muestra asociacion fuerte con \texttt{discoverymethod}; por tanto, debe interpretarse como evidencia mixta, no como estructura fisica pura.

La complejidad termica es especialmente cautelosa por su alta imputacion. La densidad funciona como control metodologico porque muchos valores de \texttt{pl\_dens} son derivados de masa y radio. La conclusion principal no es una taxonomia final, sino una clasificacion reproducible de regiones: algunas fisicamente interpretables, algunas observacionalmente sospechosas, muchas mixtas y otras debiles. Este resultado establece una base medible para una siguiente etapa de validacion, sin sobreafirmar la topologia real del catalogo.
